\chapter{PROBLEM 5}
\begin{center}\large\bfseries FAN-OUT EVENT DISTRIBUTION/PUB-SUB\end{center}

\section{Rancangan Solusi dan Mekanisme}
Pada Problem 5, event dari satu producer perlu diterima oleh beberapa consumer, termasuk consumer yang sempat berhenti atau baru ditambahkan. Rancangan menggunakan broker persisten sebagai perantara. Setiap consumer memiliki posisi baca sendiri, sehingga dapat mengejar event setelah pulih dan tidak mengambil jatah pesan consumer lain.

\subsection{Producer, stream, dan consumer}
Alur distribusi dimulai ketika Aggregator menyimpan perubahan hazard. Data kanonik, revisi, dan pesan outbox ditulis dalam satu transaksi, sehingga perubahan yang tersimpan mempunyai catatan untuk dikirim. Relay kemudian menerbitkan envelope ke subject \path{hazard.events.v1.BMKG} atau \path{hazard.events.v1.PVMBG}. Stream \texttt{HAZARD\_EVENTS} menyimpan pesan dari kedua subject melalui filter \texttt{hazard.events.v1.*}.

Producer mempublikasikan pesan berdasarkan subject dan kontrak event, sehingga penambahan pembaca tidak memerlukan perubahan pada producer. Daftar consumer serta posisi baca masing-masing dikelola oleh broker.

Dashboard dan Notifier memakai durable yang berbeda, yaitu \texttt{dashboard-v1} dan \texttt{notifier-v1}. Broker menyimpan posisi baca setiap durable meskipun aplikasinya berhenti. Dengan posisi yang terpisah, masing-masing aplikasi dapat menerima seluruh event dan memberikan ACK sendiri.

Konfigurasi berikut mengatur berapa lama event tersedia, kapan pesan dikirim ulang, dan berapa banyak pesan yang dapat diproses bersamaan. Batas retention menentukan kesempatan consumer yang pulih atau baru ditambahkan untuk membaca event lama.

\begin{table}[H]
\centering\small
\begin{tabularx}{\textwidth}{>{\raggedright\arraybackslash}p{4.8cm}Y}
\toprule
Konfigurasi & Nilai dan fungsi \\
\midrule
Storage / replica / retention & File / 1 / LimitsPolicy; ACK consumer tidak menghapus event dari stream. \\
Batas retention & 24 jam, 256 MiB, 100.000 pesan; batas yang tercapai lebih dahulu menentukan kapasitas. \\
Batas satu pesan & 64 KiB, termasuk ukuran yang diperhitungkan publisher. \\
Saat kapasitas penuh & DiscardNew; publish baru ditolak, bukan menghapus pesan lama hanya untuk menerima pesan baru. \\
Dedup broker & Message ID \texttt{hazard\_id:revision}; window 2 menit. \\
ACK / redelivery & AckExplicit, AckWait 30 detik, MaxDeliver $-1$; retry tidak dibatasi jumlahnya. \\
Pull / inflight & Batch 10, expiry fetch 2 detik, MaxAckPending 64. \\
\bottomrule
\end{tabularx}
\caption{Konfigurasi JetStream dan consumer yang digunakan pada M1.}
\end{table}

\subsection{Semantik pengiriman dan idempotensi}
Semantik pengiriman sistem adalah \textbf{at-least-once}, sehingga sebuah pesan dapat diterima lebih dari sekali. Misalnya, broker sudah menerima publish, tetapi relay gagal mencatat status published di database. Relay akan mencoba kembali. Hal serupa dapat terjadi ketika consumer sudah melakukan commit lokal, tetapi proses mati sebelum ACK diterima broker. Message ID \texttt{hazard\_id:revision} membantu broker mengenali publish ulang selama dedup window 2 menit, tetapi perlindungan tersebut memiliki batas waktu.

Untuk menangani pengiriman ulang, setiap consumer menyimpan \texttt{message\_id} pada inbox dengan constraint unik. Pencatatan inbox dan hasil olahan dilakukan dalam satu transaksi SQLite, kemudian consumer mengirim ACK. Jika pesan yang sama datang kembali, inbox menunjukkan bahwa hasilnya sudah diterapkan sehingga side effect lokal tidak diulang. Dashboard hanya menerima revisi yang lebih baru, sedangkan Notifier menggunakan identitas pesan yang unik untuk alert-nya.

Pesan dengan envelope tidak valid disimpan ke quarantine sebelum ACK, agar pesan rusak tidak terus diulang dan payload tetap tersedia untuk diperiksa. Jika penyimpanan gagal, consumer meminta pengiriman ulang sehingga pesan belum dianggap selesai.

Saat publish gagal, relay mencoba kembali dengan backoff 1 sampai 30 detik. Pending outbox dibatasi 10.000 pesan; jika batas ini tercapai, ingest tidak terus menambah antrean. Penghentian satu consumer karena itu tidak langsung menghalangi producer, selama broker masih dapat menyimpan event. Namun, kapasitas broker dan disk tetap menjadi batas: dengan DiscardNew, publish baru ditolak ketika stream penuh dan pesan tetap menunggu di outbox.

\section{Alasan Keputusan dan Trade-off}
Pemilihan mekanisme mempertimbangkan dua risiko: perubahan data tidak sempat dipublikasikan, dan consumer tidak tersedia saat event terbit. Outbox menangani kebutuhan publish yang persisten, sedangkan stream dan durable consumer menangani penyimpanan serta pembacaan event.

\subsection{Outbox dan fan-out}
Outbox dipilih untuk menangani celah antara penyimpanan data dan publish. Jika publish langsung dilakukan setelah commit, proses dapat mati ketika data sudah tersimpan tetapi pesan belum terkirim. Sebaliknya, publish sebelum commit dapat mengirim event dari transaksi yang kemudian gagal. Outbox mencatat kebutuhan publish dalam transaksi data yang sama, lalu relay mencobanya kembali. Konsekuensinya, ada jeda sebelum event sampai ke consumer dan kemungkinan publish berulang.

Jika kedua aplikasi berbagi satu durable, pesan akan dibagi sebagai pekerjaan untuk pembaca durable tersebut. Pola itu tidak memenuhi kebutuhan fan-out, yaitu event yang sama dapat diterima oleh setiap aplikasi. Karena itu, kedua consumer memakai pull consumer yang independen meskipun stream-nya sama.

\subsection{Broker, pola baca, dan retention}
Pull dipilih agar consumer meminta batch berikutnya setelah pengolahan batch sebelumnya selesai. Dibanding push, laju pengambilan lebih langsung dikendalikan aplikasi; ukuran batch dan batas ACK membatasi pesan yang sedang ditangani. Saat stream kosong, consumer tetap melakukan fetch dengan timeout.

LimitsPolicy mempertahankan event sampai batas retention, walaupun consumer telah memberikan ACK. Dengan DeliverAll, consumer baru dapat membaca event lama yang masih tersimpan. InterestPolicy juga mendukung fan-out, tetapi menghapus pesan setelah seluruh consumer yang berkepentingan memberikan ACK. Sementara itu, WorkQueuePolicy ditujukan untuk pembagian pekerjaan dengan satu pemroses per pesan [8]. LimitsPolicy dipilih agar penambahan consumer pada P5 dapat disertai replay, dengan biaya penyimpanan event selama masa retention.

Alternatif producer memanggil setiap consumer melalui HTTP membutuhkan daftar tujuan dan retry per tujuan di producer. NATS Core memberi fan-out sederhana, tetapi tidak menyimpan backlog durable untuk consumer yang mati. JetStream dipilih agar replay dan posisi baca dikelola broker. Kafka dan RabbitMQ juga dapat memenuhi kebutuhan ini; setup JetStream satu node dipilih karena fungsi stream dan durable sudah tersedia dalam executable NATS. Pertimbangan teknologi broker dijelaskan pada Bab 3.

\section{Batasan yang Disadari}
Jaminan idempotensi berlaku pada hasil olahan di database lokal. Seluruh alur belum dijamin exactly-once; panggilan ke layanan SMS eksternal, jika ditambahkan, berada di luar transaksi SQLite dan memerlukan penanganan idempotensi sendiri.

Consumer dapat mengejar backlog selama event yang dibutuhkannya masih tersedia di stream. Durable mempertahankan posisi baca, tetapi tidak memperpanjang retention. Jika consumer berhenti melewati batas retention, event yang sudah kedaluwarsa tidak dapat diambil kembali dari broker.

Konfigurasi memakai satu replica untuk deployment pada satu host, sehingga belum menyediakan failover antarmesin. Di sisi consumer, idempotensi bergantung pada riwayat inbox yang tersimpan di volume lokal. Jika volume tersebut dihapus, consumer kehilangan catatan pesan yang pernah diolah. Backup serta pemulihan volume belum menjadi bagian dari pengujian M1. Retention broker dan keberlangsungan volume karena itu menjadi bagian dari batas jaminan pengiriman pada M1.

\section{Bukti Penyelesaian}
Pengujian P5 mengikuti posisi event sebelum, selama, dan setelah Dashboard berhenti. Setelah recovery diperiksa, pengujian menambahkan consumer ketiga untuk membuktikan producer tidak perlu diubah. Buktinya ditunjukkan melalui command berikut dan hasil trial integrasi [10].

\subsection{Mencatat posisi dua durable consumer}
Posisi awal dicatat agar event yang terbit selama gangguan dapat dibedakan dari event setelah recovery. Pengujian membaca endpoint monitoring JetStream melalui \texttt{wget} di dalam container NATS karena port monitoring tidak dibuka ke host. Query \texttt{streams=true\&consumers=true} meminta informasi stream beserta consumer-nya, lalu \texttt{jq} memilih dua durable utama.

Pada output, \texttt{last\_seq} adalah sequence terakhir stream, sedangkan \texttt{ack} menampilkan ACK floor consumer dalam sequence stream. ACK floor menunjukkan posisi yang sudah diakui secara berurutan. \texttt{pending} menghitung pesan yang belum dikirim, dan \texttt{inflight} menghitung pesan yang sudah dikirim tetapi belum di-ACK. Karena dua jenis antrean ini berbeda, pending nol saja belum membuktikan recovery selesai.
\terminalfigure{p5-before}{Posisi awal stream dan dua durable utama sebelum Dashboard dihentikan.}{fig:p5-before}

Pada snapshot awal, Dashboard dan Notifier sama-sama memiliki ACK 308, sesuai sequence terakhir stream. Kesamaan posisi tersebut menunjukkan bahwa kedua consumer telah mengakui event hingga sequence terakhir stream.

\subsection{Menghentikan satu consumer dan membaca backlog}
Dashboard dihentikan melalui Compose, sementara producer dan Notifier dibiarkan berjalan. Jeda 15 detik memberi kesempatan bagi mock untuk menghasilkan event baru dan Aggregator untuk meneruskannya ke broker.
\terminalfigure{p5-stop}{Menghentikan Dashboard tanpa menghentikan producer, broker, atau Notifier.}{fig:p5-stop}

Command monitoring yang sama kemudian dijalankan ulang. Sequence stream naik dari 308 menjadi 311 dan ACK Notifier ikut mencapai 311. ACK Dashboard tetap 308 dengan tiga pesan pending. Jadi, event baru tetap tersimpan dan diolah Notifier saat Dashboard tidak aktif. Sequence 311 dicatat sebagai target yang harus dicapai Dashboard setelah kembali berjalan.
\terminalfigure{p5-stopped}{Stream dan Notifier terus maju, sedangkan backlog Dashboard bertambah.}{fig:p5-stopped}

Dashboard dijalankan kembali, kemudian pengujian menunggu 35 detik sebelum membaca posisi ACK. Jeda ini juga memberi ruang untuk redelivery jika masih ada pesan yang dikirim sebelum consumer berhenti; konfigurasi AckWait adalah 30 detik.
\terminalfigure{p5-recovery-control}{Menjalankan kembali Dashboard dan menunggu pemrosesan backlog.}{fig:p5-recovery-control}

Pada pemeriksaan berikutnya, ACK Dashboard mencapai 318, melewati target 311 yang dicatat saat ia mati. Notifier juga mencapai 318, dan kedua consumer tidak mempunyai pesan pending maupun inflight pada snapshot tersebut. Hasil ini membuktikan Dashboard memproses event yang terbit selama gangguan, bukan hanya menerima event setelah startup.
\terminalfigure{p5-recovery}{Dashboard sudah mengakui seluruh pesan hingga melewati target periode gangguan.}{fig:p5-recovery}

Inspeksi database lokal kemudian menunjukkan bahwa kedua consumer masing-masing memiliki 319 pesan pada inbox dan quarantine nol. Jumlahnya dapat lebih tinggi dari snapshot broker sebelumnya karena stream terus menerima event selama inspeksi. Inbox yang terpisah menunjukkan bahwa setiap aplikasi mencatat pengolahannya secara independen.
\terminalfigure{p5-inbox}{Hasil olahan berada pada inbox lokal masing-masing consumer.}{fig:p5-inbox}

\subsection{Menambahkan consumer tanpa mengubah producer}
Penambahan consumer diuji pada sesi terpisah. Sebelum Third Consumer dijalankan, ID, waktu mulai, dan restart count Aggregator disimpan. Profile \texttt{demo} kemudian menambahkan container Third Consumer dengan durable \texttt{third-consumer-v1}, tanpa membangun atau menjalankan ulang Aggregator. Consumer ini memakai binary Notifier dengan durable dan volume miliknya sendiri.
\terminalfigure{p5-third-start}{Cuplikan awal dan akhir command startup Third Consumer melalui profil demo; bagian tengah log build tidak ditampilkan.}{fig:p5-third-start}

Pada sesi tersebut, snapshot monitoring menunjukkan durable baru telah mencapai ACK 192, sama dengan posisi stream dan kedua durable lainnya. Karena consumer baru memakai DeliverAll, pesan yang masih tersimpan sebelum ia dibuat juga tersedia untuk dibaca; producer tidak perlu mengirim ulang secara khusus kepadanya.
\terminalfigure{p5-third}{Durable baru membaca stream yang sama sampai posisi ACK yang ditampilkan broker.}{fig:p5-third}

Perbandingan kondisi Aggregator sebelum dan sesudah menghasilkan file yang identik. Inspeksi Third Consumer juga menunjukkan inbox 193 dengan quarantine nol. Jumlah inbox diambil setelah snapshot broker, ketika stream masih menerima event. Dengan demikian, penambahan consumer berhasil tanpa perubahan container producer, dan consumer baru berhasil menyimpan hasil pengolahannya.
\terminalfigure{p5-third-producer}{Producer tidak berubah dan consumer tambahan memiliki inbox sendiri.}{fig:p5-third-producer}

Hubungan hasil pengujian dengan kriteria P5 dirangkum pada tabel berikut.
\begin{table}[H]
\centering\small
\begin{tabularx}{\textwidth}{>{\raggedright\arraybackslash}p{2cm}Y}
\toprule
Kriteria & Bukti \\
\midrule
P5.1 & Publish relay ke subject broker pada kode dan kenaikan sequence ketika Dashboard mati. \\
P5.2 & Posisi durable terpisah pada Gambar~\ref{fig:p5-before} dan inbox masing-masing pada Gambar~\ref{fig:p5-inbox}. \\
P5.3 & Sequence 311 sebagai target gangguan, lalu ACK Dashboard 318 pada Gambar~\ref{fig:p5-stopped} dan \ref{fig:p5-recovery}. \\
P5.4 & Startup Third Consumer, posisi ACK, dan kondisi producer pada Gambar~\ref{fig:p5-third-start}--\ref{fig:p5-third-producer}. \\
P5.5 & Broker, pola fan-out, semantik at-least-once, dan konsekuensi idempotensi consumer pada Bagian~\ref{sec:p5-di-laporan}. \\
\bottomrule
\end{tabularx}
\end{table}

\section{Jawaban atas Butir ``Di laporan''}
\label{sec:p5-di-laporan}
Broker yang dipilih adalah NATS JetStream dengan pola fan-out melalui stream \texttt{HAZARD\_EVENTS}. Setiap aplikasi memakai durable pull consumer sendiri, sehingga seluruh aplikasi dapat membaca event yang sama tanpa daftar tujuan pada producer. LimitsPolicy dan DeliverAll memungkinkan consumer yang pulih atau baru ditambahkan membaca event yang masih berada dalam retention.

Semantik pengiriman adalah \textbf{at-least-once}: kegagalan pencatatan status publish atau ACK dapat menyebabkan pesan dikirim ulang. Karena itu, consumer harus idempoten. Kelompok menerapkannya dengan unique \texttt{message\_id} pada inbox dan transaksi SQLite yang mencatat inbox bersama hasil olahan sebelum ACK. Dedup broker membantu selama window 2 menit, tetapi inbox tetap diperlukan untuk redelivery di luar window tersebut. Jaminan ini berlaku pada hasil lokal; panggilan notifikasi eksternal kelak memerlukan penanganan idempotensi tersendiri.
